\chapter{散射理论}
\label{ch:scattering}

\section{散射截面}
上一章讨论了如何从量子态预测仪器记录各个结果的概率。散射实验测量粒子进入给定角窗的计数，它还取决于入射通量和束流照射到的靶粒子数。用这两个量归一化计数，就得到可在不同装置间比较的散射截面。截面的角分布又能提供哪些关于相互作用的信息？

\begin{definition}{微分散射截面}{qm-cross-section}
    对单个靶散射中心，进入给定末态通道的微分截面(differential cross section)定义为
    \begin{equation}
        \frac{\d\sigma}{\nd\Omega}
        =\frac{\text{单位时间进入单位立体角的粒子数}}{\text{入射通量密度}}.
    \end{equation}
    对探测器接收的立体角积分得到相应的部分截面，量纲为面积。
\end{definition}
设入射通量密度为$j_{\incoming}$，量纲为面积$^{-1}$时间$^{-1}$；
束流照到$N_T$个彼此独立的有效散射中心，
探测器接收$\Delta\Omega$立体角，效率为$\eta$，计数时间为$\tau$。
若角窗内微分截面近似常量，单次散射给出
\begin{equation}
    N_\mathrm{det}-N_\mathrm{bg}
    \simeq \eta\tau j_{\incoming} N_T
        \frac{\d\sigma}{\nd\Omega}\Delta\Omega.
\end{equation}
因此从实验计数提取截面，要分别标定通量、靶厚度、效率和角接受度（图\ref{fig:scattering-aperture}）。
\begin{figure}[htbp]
    \centering\tikzfigure[width=.56\linewidth]{scattering-aperture}
    \caption{散射计数的几何。$\theta$ 是探测器中心相对入射束的散射角；阴影为接收锥的截面，从靶心看去，整个探测面张开的立体角是 $\Delta\Omega$。}\label{fig:scattering-aperture}
\end{figure}
当靶太厚时，后续再次散射改变出射角分布，
计数不再与$N_T$简单成正比。

改变入射动量可以探测相互作用的不同空间尺度。测量散射后的能量，还可区分靶仍在原内部态的弹性过程与激发、反应等非弹性过程。
有限探测器实际积分一个角度和能量窗口；多次散射、靶的热运动及非弹性通道(inelastic channel)要另行建模。
本章先取非相对论、无自旋、短程势(short-range potential)的弹性散射(elastic scattering)，$\mu$为约化质量。

自由能量为$E_k=\hbar^2k^2/(2\mu)$。采用三维波矢态约定
\begin{equation}
    \brkt{\bm k}{\bm k'}=\delta^{(3)}(\bm k-\bm k'),\qquad
    \brkt{\bm x}{\bm k}=(2\pi)^{-3/2}\e^{\i\bm k\cdot\bm x}.
\end{equation}
$\ket{\bm k}$的动量是$\hbar\bm k$；本章的连续归一化与按动量归一化相差$\hbar$因子。
入射实验应理解为窄而可归一化的动量分布，精确$\ket{\bm k}$是用于取散射极限的广义态。

\subsection{探测面上的球面波}
取入射方向为$z$轴，探测器到靶心的距离为$r$。在远大于势的作用范围处，
入射波的波面近似平面；从靶向某一方向散出的波穿过的面积却随$r^2$增长。
为使每秒穿过相同立体角的粒子数与$r$无关，出射波的振幅必须按$1/r$衰减。
采用时间因子$\exp(-\i E_k t/\hbar)$，于是可把散射态的位置分量写为
\begin{equation}
    \brkt{\bm x}{\bm k^{(+)}}\ \sim\
    C\Bigb{\e^{\i kz}+f(\bm k',\bm k)\frac{\e^{\i kr}}{r}},
    \qquad C=(2\pi)^{-3/2},\quad \bm k'=k\,\bm x/r .
    \label{eq:outgoing-wave-pattern}
\end{equation}
相位$kr-E_kt/\hbar$恒定的波面向外移动，$f$就是散射振幅(scattering amplitude)，其量纲是长度。
在直射束以外，探测器只接收出射波。将出射分量$\psi_{\mathrm{sc}}$代入自由粒子的径向概率流
$j_r=(\hbar/\mu)\Im(\psi_{\mathrm{sc}}\conj\p_r\psi_{\mathrm{sc}})$，得
\begin{equation}
    j_r=\frac{\hbar k}{\mu}\abs{C}^2\frac{\abs{f}^2}{r^2},
    \qquad
    \frac{j_r r^2\d\Omega}{j_{\incoming}}=\abs{f}^2\d\Omega,
    \qquad j_{\incoming}=\frac{\hbar k}{\mu}\abs{C}^2 .
\end{equation}
分子是穿过探测面的粒子数率，分母是照到靶上的通量密度。
这就把实验定义的$\d\sigma/\nd\Omega$与理论中的$\abs{f}^2$联系起来。
在前向还要考虑入射波与出射波的干涉，后面的光学定理将给出它与总截面的关系。

\section{散射态、预解式与跃迁算符}

远处探测器的角分布给出散射振幅的模平方。给定靶的势能算符 $\op{V}$，怎样求出这个振幅？
把定态方程写成 $(E-\op{H}_0)\ket{\psi}=\op{V}\ket{\psi}$，就能把右边看成分布在靶区内的波源。
需要找到将源传播到探测器的算符，并选对边界条件，使靶激起的波向外传播。

\begin{definition}{预解式与出射边界值}{qm-resolvent}
    对自伴Hamilton算符$\op{H}$，复数$z$不在其谱上时，
    \begin{equation}
        \op{G}(z)=(z\one-\op{H})^{-1}
    \end{equation}
    称为预解式(resolvent)。后文的 $z-\op{H}$ 是 $z\one-\op{H}$ 的简写。
    预解式把给定的源态$\ket{\eta}$映到
    满足$(z-\op{H})\ket{\psi}=\ket{\eta}$的态。
    自由系统的预解式记为$\op{G}_0(z)=(z-\op{H}_0)^{-1}$。
    连续谱上的$\op{G}_0^{(+)}(E)$是$z=E+\i\eta$、
    $\eta\to0^+$的边界值。
\end{definition}

实验中的入射态先在远离散射中心处自由传播。
在存在相应极限的散射子空间中，定义入、出态映射
\begin{subequations}
    \begin{alignat}{1}
        \op{\Omega}_{\incoming} &=\lim_{t\to-\infty}\e^{\i \op{H}t/\hbar}\e^{-\i \op{H}_0t/\hbar},\\
        \op{\Omega}_{\outgoing} &=\lim_{t\to+\infty}\e^{\i \op{H}t/\hbar}\e^{-\i \op{H}_0t/\hbar}.
    \end{alignat}
\end{subequations}
$\op{\Omega}_{\incoming}$把自由入射态映到相互作用中的态，
$\op{S}=\op{\Omega}_{\outgoing}\adjoint\op{\Omega}_{\incoming}$则将入射自由标签映到最终出射标签。
$\ket{\bm k^{(+)}}$表示由$\ket{\bm k}$指定的入射散射态(scattering state)。

令$\op{H}=\op{H}_0+\op{V}$，定义自由预解式
$\op{G}_0^{(+)}(E)=(E-\op{H}_0+\i0)^{-1}$。
Lippmann--Schwinger方程写作
\begin{equation}
    \ket{\bm k^{(+)}}=\ket{\bm k}+\op{G}_0^{(+)}(E_k)\op{V}\ket{\bm k^{(+)}}.
\end{equation}
它由$(E-\op{H}_0)\ket{\bm k^{(+)}}=\op{V}\ket{\bm k^{(+)}}$配合出射边界条件得到；
$+\i0$表示先取$E+\i\eta$、$\eta>0$，再令$\eta\to0^+$。
可从位置表象看出这个符号的作用。记 $R=\abs{\bm x-\bm x'}$，自由预解式在位置表象中的核是一个复值函数，
\begin{equation}
    G_0^{(+)}(\bm x,\bm x';E_k)=\mel{\bm x}{\op{G}_0^{(+)}(E_k)}{\bm x'}
    =-\frac{\mu}{2\pi\hbar^2}\frac{\e^{\i kR}}{R}.
\end{equation}
当能量有正虚部时，取正实部的波数也有正虚部，$\e^{\i kR}$随距离衰减；
令虚部趋零留下向外传播的波。为确定核的系数，取包围波源的小球，
将 $(E_k+\hbar^2\nabla_{\bm x}^2/2\mu)G_0^{(+)}(\bm x,\bm x';E_k)=\delta^{(3)}(\bm x-\bm x')$ 在球内积分，
球缩小时只有$1/R$项的梯度通量$-4\pi$留下，恰好给出右端的1。
\begin{definition}{跃迁算符}{qm-scattering-t}
    跃迁算符(transition operator)$\op{T}(E)$把给定的自由入射态映到激起散射波的源态，满足
    \begin{equation}
        \op{T}(E)=\op{V}+\op{V}\op{G}_0^{(+)}(E)\op{T}(E),\qquad
        \op{T}(E_k)\ket{\bm k}=\op{V}\ket{\bm k^{(+)}}.
    \end{equation}
\end{definition}
迭代为$\op{T}=\op{V}+\op{V}\op{G}_0\op{V}+\op{V}\op{G}_0\op{V}\op{G}_0\op{V}+\cdots$，一级 Born 近似取 $\op{T}\simeq \op{V}$。

同一个结构还可由预解式恒等式看出：
\begin{equation}
    \op{G}=\op{G}_0+\op{G}_0\op{V}\op{G},\qquad \op{T}=\op{V}+\op{V}\op{G}\op{V}.
\end{equation}
第一式由$(z-\op{H})=(z-\op{H}_0)-\op{V}$左右乘相应逆算符得到；
将它代入第二式，即得$\op{T}=\op{V}+\op{V}\op{G}_0\op{T}$。
算符$\op{T}$与$\op{V}$同样具有能量量纲；
在本章的三维$\delta$归一化下，其矩阵元的量纲为能量乘长度的三次方。
壳上(on-shell)矩阵元连接能量相同的入射与出射通道。

截断的精度由后续各阶在所研究态上的贡献决定。
接近束缚态阈值或共振时，微弱势也可能产生强散射。

要从$\op{T}$读出$f$，取探测器远离靶，记$\bm n=\bm x/r$。
对靶区内的$\bm x'$，有$R=r-\bm n\cdot\bm x'+\bigo(1/r)$。
于是核中的$\e^{\i kR}/R$在远处成为$(\e^{\i kr}/r)\e^{-\i\bm k'\cdot\bm x'}$。
与式\eqref{eq:outgoing-wave-pattern}比较，球面波的系数满足
\begin{equation}
    C f=-\frac{\mu}{2\pi\hbar^2}\int\d^3 x'\,
        \e^{-\i\bm k'\cdot\bm x'}\mel{\bm x'}{\op{V}}{\bm k^{(+)}}
    =-\frac{\mu}{2\pi\hbar^2 C}\mel{\bm k'}{\op{T}(E_k)}{\bm k}.
\end{equation}
最后一步使用$\brkt{\bm k'}{\bm x'}=C\e^{-\i\bm k'\cdot\bm x'}$。
由连接入、出自由态的散射算符 $\op{S}$，可将这些关系写成
\begin{equation}
    \begin{alignedat}{1}
        \mel{\bm k'}{\op{S}}{\bm k} &=\delta^{(3)}(\bm k'-\bm k)-2\pi\i\delta(E_{k'}-E_k)\mel{\bm k'}{\op{T}(E_k)}{\bm k},\\
        f(\bm k',\bm k)&=-\frac{4\pi^2\mu}{\hbar^2}\mel{\bm k'}{\op{T}(E_k)}{\bm k},\qquad
            \frac{\d\sigma}{\nd\Omega}=\abs{f}^2.
        \label{eq:scatnorm}
    \end{alignedat}
\end{equation}
最后一式限同一弹性通道；不同通道通常还带末初速度比。
完整的散射子空间中，概率守恒要求$\op{S}\adjoint \op{S}=\one$。
对短程自伴势，合适的入出态极限存在时可构造这些算符；Coulomb长程尾需要修改渐近态。

\section{Born近似}
在跃迁算符的展开中，第一项只包含一次势的作用。如果反复散射的振幅足够小，
就能直接用势的空间形状预言角分布。以下把这一近似与入射通量和末态数联系起来，
由此确定各散射角对应的势场空间频率。

为核对式\eqref{eq:scatnorm}的因子，暂用体积$\mathcal V$中单位归一的动量态。
对局域势 $\op{V}=V(\hatbm{X})$，动量态之间的矩阵元是$\widetilde V(\bm q)/\mathcal V$，其中
\begin{equation}
    \widetilde V(\bm q)=\int\d^3x\e^{-\i\bm q\cdot\bm x}V(\bm x),\qquad
    \bm q=\bm k'-\bm k.
\end{equation}
末态数为 $\mathcal V k'^2\d k'\d\Omega/(2\pi)^3$。利用 $E_{k'}=\hbar^2 k'^2/(2\mu)$，
能量壳积分只选中正根 $k'=k$，给出
\[
    \intRplus k'^2\delta(E_{k'}-E_k)\d k'
    =\frac{k^2}{\hbar^2 k/\mu}=\frac{\mu k}{\hbar^2}.
\]
把这些因子按实验定义依次组合，单位立体角内的跃迁率为
\begin{equation}
    \frac{\d\Gamma}{\nd\Omega}
    =\frac{2\pi}{\hbar}
        \frac{\abs{\widetilde V(\bm q)}^2}{\mathcal V^2}
        \frac{\mathcal V}{(2\pi)^3}
        \frac{\mu k}{\hbar^2}.
\end{equation}
第一项是黄金规则系数，第二项是盒归一化的矩阵元平方，
第三项计数末态，最后一项来自能量$\delta$函数的Jacobian。
单位归一化入射态在盒内的粒子数密度为$1/\mathcal V$，
速度为$\hbar k/\mu$，故通量密度为$j_{\incoming}=\hbar k/(\mu\mathcal V)$。
将上式除以这个通量，得到截面的Born近似。振幅的符号由前面选定的出射波约定给出。
\begin{equation}
    \Bigp{\frac{\d\sigma}{\nd\Omega}}_{\Born}
    =\frac{\mu^2}{4\pi^2\hbar^4}\abs{\widetilde V(\bm q)}^2,\qquad
    f_{\Born}=-\frac\mu{2\pi\hbar^2}\widetilde V(\bm q).
\end{equation}
所有任意量子化体积均消去，$f$有长度量纲。

盒归一化态与$\delta$归一化态的关系为
$\ket{\bm k}_{\boxnorm}=(2\pi)^{3/2}\ket{\bm k}/\sqrt{\mathcal V}$。
因此对局域势，
\begin{equation}
    \mel{\bm k'}{\op{V}}{\bm k}=\frac{\widetilde V(\bm k'-\bm k)}{(2\pi)^3}.
\end{equation}
代入式\eqref{eq:scatnorm}中$f$与$\op{T}$的关系，
在$\op{T}\simeq \op{V}$时恰好得到所列$f_{\Born}$。
这使连续态约定与通量计算使用同一归一化。

取短程Yukawa势$V(r)=g\e^{-r/a}/r$，$a>0$，$g$量纲为能量乘长度。
把积分的极轴选为$\bm q$，以$u=\cos\vartheta$为积分变量，
先做方位角和极角积分：
\begin{equation}
    \int\d\Omega_{\bm x}\e^{-\i qr\cos\vartheta}
    =2\pi\int_{-1}^{1}\e^{-\i qru}\d u
    =4\pi\frac{\sin(qr)}{qr}.
\end{equation}
$r^2$的体积因子与势中的$1/r$、角积分中的$1/r$相消，留下
\begin{equation}
    \widetilde V(q)=\frac{4\pi g}{q}\intRplus\e^{-r/a}\sin(qr)\d r
    =\frac{4\pi g}{q^2+a^{-2}},\qquad q=2k\sin\frac\theta2.
\end{equation}
因而$f_{\Born}=-2\mu g/[\hbar^2(q^2+a^{-2})]$。
低能弱耦合的自然检验参数为$\mu\abs{g}a/\hbar^2\ll1$；更一般应检查Born后项。
取 $a\to\infty$ 可得未屏蔽势的 Fourier 变换（取 $q>0$）。构造散射态时还要使用适合 Coulomb 长程尾的渐近边界条件。

\section{分波、幺正性与低能共振}
势强到需要保留多次散射时，一级Born结果可能失准。若靶具有球对称性，
角动量守恒使不同角动量的分量独立传播。分别求出势对各分量的作用，再将它们的振幅相加，就能得到探测器所见的角分布。

\begin{definition}{分波与散射相移}{qm-phase-shift}
    对无自旋粒子在实短程中心势中的弹性散射，按轨道角动量 $\ell$ 分出的成分称为分波(partial wave)。
    固定波数 $k>0$，在势外远处，自由径向驻波乘以 $r$ 后正比于 $\sin(kr-\ell\pi/2)$。
    相互作用将其变为正比于 $\sin(kr-\ell\pi/2+\delta_\ell)$ 的驻波；$\delta_\ell$ 称为相移(phase shift)，模 $\pi$ 确定。
\end{definition}
令$\xi=kr-\ell\pi/2$，把驻波拆成向外与向内的两个指数项，
\begin{equation}
    \sin(\xi+\delta_\ell)
    =\frac{\e^{\i\delta_\ell}\e^{\i\xi}
         -\e^{-\i\delta_\ell}\e^{-\i\xi}}{2\i}.
\end{equation}
与自由驻波比较，出射分量多出的相位是$\delta_\ell$，入射分量的是$-\delta_\ell$。
将入射分量固定为自由波的值，出射分量便乘上 $S_\ell=\e^{2\i\delta_\ell}$。这是散射算符 $\op{S}$ 在该弹性分波上的本征值，指数中的因子 2 来自出入两支的相对相位。
把入出方向态分解到角动量基，利用角度基的加法定理，得到
\begin{equation}
    f(\theta)=\frac1{2\i k}\sum_{\ell=0}^\infty(2\ell+1)(S_\ell-1)P_\ell(\cos\theta)
    =\frac1k\sum_\ell(2\ell+1)\e^{\i\delta_\ell}\sin\delta_\ell P_\ell(\cos\theta).
\end{equation}
$P_\ell$是$\ell$阶Legendre多项式，来自角动量方向基的加法定理。
利用正交关系
\[
    \int\d\Omega P_\ell(\cos\theta)P_{\ell'}(\cos\theta)=\frac{4\pi}{2\ell+1}\delta_{\ell\ell'}
\]
得到
\begin{equation}
    \sigma_{\elastic}=\frac{4\pi}{k^2}\sum_\ell(2\ell+1)\sin^2\delta_\ell .
\end{equation}
\begin{theorem}{光学定理}{qm-optical}
    对短程相互作用，若所有开放通道的总通量守恒，本章的出射波约定给出
    \begin{equation}
        \sigma_\tot=\frac{4\pi}{k}\Im f(0).
    \end{equation}
    左边计入所有末态通道，右边取入射通道的前向弹性散射振幅。
\end{theorem}
\begin{proof}[纯弹性中心势下的证明]
    在前向$\theta=0$，$P_\ell(1)=1$，并且
    $\Im(\e^{\i\delta_\ell}\sin\delta_\ell)=\sin^2\delta_\ell$。
    所以
    \begin{equation}
        \Im f(0)=\frac1k\sum_\ell(2\ell+1)\sin^2\delta_\ell
        =\frac{k}{4\pi}\sigma_{\elastic}.
    \end{equation}
    角积分中不同$\ell$的交叉项因Legendre正交性消失；
    前向虚部与总计数的关系因此直接来自各弹性通道的幺正性。
\end{proof}

有开放非弹性通道时，弹性块可写$S_\ell=\eta_\ell\e^{2\i\delta_\ell}$，$0\leq\eta_\ell\leq1$，
光学定理(optical theorem)中的总截面(total cross section)必须包括这些通道。

\begin{example}{从硬球边界读出相移}{qm-hard-sphere-phase}
    半径为$a$的硬球要求波在$r=a$处为零。考虑角度上均匀的$\ell=0$分波（$s$波），
    写$u_0=rR_0$，其中$R_0$是态的径向分量。
    球外的自由方程为$u_0''+k^2u_0=0$，满足硬壁条件的解正比于$\sin[k(r-a)]$。
    与无势时的$\sin(kr)$对照，硬壁使球外的波节向外移了$a$（图\ref{fig:hard-sphere-nodes}），故$\delta_0=-ka$（模$\pi$）。
    当$ka\ll1$，其余分波贡献很小，$f\simeq\e^{\i\delta_0}\sin\delta_0/k\simeq-a$。
    探测器转到各个方向都测得约$a^2$的微分截面；对整个球面积分得$\sigma\simeq4\pi a^2$。
    这说明波长远大于靶时，按出射通量定义的截面可以超过硬球的几何投影面积$\pi a^2$。
\end{example}

\begin{figure}[htbp]
    \centering
    \tikzfigure{hard-sphere-nodes}
    \caption{取 $ka=\pi/4$，纵轴为 $u_0=rR_0$，两条曲线按相同峰值绘制。灰区为硬球内部；比较虚线与实线在 $r/a=4$ 和 $5$ 处的过零点，即可读出波节位移。}
    \label{fig:hard-sphere-nodes}
\end{figure}

令$a$为势的有效作用范围。对短程势，$ka\ll1$且不处于高分波阈值共振时，通常$s$波占优。
有效力程(effective range)展开为
\begin{equation}
    k\cot\delta_0=-\frac1{a_s}+\frac12r_e k^2+\cdots,
    \qquad f_0=\frac1{k\cot\delta_0-\i k}.
\end{equation}
忽略$r_e$项得到$\sigma\simeq4\pi a_s^2/(1+k^2a_s^2)$。
$a_s$是散射长度(scattering length)，$r_e$是有效力程，两者都有长度量纲。
$a_s$也能从低能径向波认出。有限$a_s$时取$\delta_0$趋零模$\pi$的分支，
$k\cot\delta_0\to-1/a_s$意味着$\delta_0\simeq-ka_s$。
在势外且$kr\ll1$时，$\sin(kr+\delta_0)\simeq k(r-a_s)$。
把这一段近似直线延长到零，其$r$轴截距就是$a_s$；对刚才的硬球，截距恰好为半径$a$。
$a_s$可正可负，由低能散射决定；当$\abs{a_s}$很大，截面受$4\pi/k^2$幺正上限约束。
这比把零能结果$4\pi a_s^2$无限外推更可靠。

\section{一维接触相互作用}
分波分析表明，低能散射对近阈束缚态很敏感。一维接触势允许精确求和多次散射振幅，可借此考察同一跃迁算符怎样联系散射与束缚态。粒子的出射方向只有反射和透射两种，两者的通量之和还可用于检验概率守恒。

取$\op{H}=\op{P}^2/(2m)+g\delta(\op{X})$，$g$有能量乘长度量纲。
这是标准一维接触相互作用(contact interaction)的自伴实现；在三维，点相互作用需要正则化并用物理散射参数确定耦合。
采用$\brkt{p}{p'}=\delta(p-p')$，势矩阵元为$g/(2\pi\hbar)$，
因此$\op{T}$的核仍可分离：$\mel{p'}{\op{T}}{p}=\tau(E)/(2\pi\hbar)$。
代入 $\op{T}=\op{V}+\op{V}\op{G}_0\op{T}$，用标量函数 $I(E)$ 记中间动量的积分，得到
\begin{equation}
    \tau(E)=\frac{g}{1-g I(E)},\qquad
    I(E)=\intR\frac{\d p}{2\pi\hbar}
    \frac1{E-p^2/(2m)+\i0}.
\end{equation}
对于$E=\hbar^2k^2/(2m)>0$，留数或主值加$\delta$分解给出$I=-\i m/(\hbar^2k)$。

展开这个积分时，使用
$(x+\i0)^{-1}=\PV(1/x)-\i\pi\delta(x)$，
其中 $\PV$ 表示 Cauchy 主值。
正能量下主值部分为零，而
\begin{equation}
    \delta\Bigp{E-\frac{p^2}{2m}}
    =\frac{m}{\hbar k}
    \bigb{\delta(p-\hbar k)+\delta(p+\hbar k)}.
\end{equation}
两个动量根都要计入，故
$I=-\i\pi(2m/\hbar k)/(2\pi\hbar)$。

定义$\alpha=\pm$表示右行或左行方向，
并以$E$和$\alpha$标记单位能量归一化的自由态。
正能量的壳上$\delta$函数带来上面的速度因子$m/(\hbar k)$，
从而散射算符在方向通道基中的矩阵元为
\begin{equation}
    S_{\alpha\beta}(E)=\delta_{\alpha\beta}
    -\frac{\i m\tau(E)}{\hbar^2k}.
\end{equation}
右行入射的透射振幅(transmission amplitude)是$t=S_{++}$，反射振幅(reflection amplitude)是$r=S_{-+}$。

在两条单位通量方向通道中，
\begin{equation}
    t=1-\frac{\i m\tau}{\hbar^2k}=\frac1{1+\i\eta},\qquad
    r=-\frac{\i m\tau}{\hbar^2k}=\frac{-\i\eta}{1+\i\eta},\quad
    \eta=\frac{mg}{\hbar^2k}.
\end{equation}
于是$\mathcal T=1/(1+\eta^2)$，$\mathcal R=\eta^2/(1+\eta^2)$，二者和为1。
在负能量$E=-\hbar^2\kappa_b^2/(2m)$时，$I=-m/(\hbar^2\kappa_b)$。

这里$\kappa_b>0$是束缚态的逆长度尺度。
此时积分没有实轴极点(pole)，直接得到
\begin{equation}
    I=-\frac{2m}{2\pi\hbar}
    \intR\frac{\d p}{p^2+\hbar^2\kappa_b^2}
    =-\frac{m}{\hbar^2\kappa_b}.
\end{equation}

极点条件$1-gI=0$仅在$g<0$时有解，
\begin{equation}
    \kappa_b=-\frac{mg}{\hbar^2},\qquad E_b=-\frac{mg^2}{2\hbar^2}.
\end{equation}
反射、透射和束缚能由同一个$\op{T}(E)$的连续谱边界值与极点给出。

\checkpoint{实势的一阶Born前向振幅为实数，
    为什么仍能满足光学定理？}{
    振幅与截面的微扰阶数不同。
    给势能算符乘上无量纲参数 $\lambda$ 以标记阶数，展开为
    \[
        \op{T}=\lambda\op{V}+\lambda^2\op{V}\op{G}_0^{(+)}\op{V}+\cdots .
    \]
    一级振幅与 $\lambda$ 成正比，其模平方给出二阶截面。下式取 $\lambda=1$。
    要在二阶检验光学定理，就必须保留二阶振幅的虚部。

    在自由能量基中，用
    $(x+\i0)^{-1}=\PV(1/x)-\i\pi\delta(x)$，
    得到前向二阶矩阵元的虚部
    \begin{equation}
        \Im\mel{\bm k}{\op{V}\op{G}_0^{(+)}(E_k)\op{V}}{\bm k}
        =-\pi\int\d^3k'\,
            \abs{\mel{\bm k'}{\op{V}}{\bm k}}^2\delta(E_k-E_{k'}).
    \end{equation}
    势为自伴保证两侧矩阵元互为共轭。
    对 $k'$ 作能量壳积分得到 $\mu k/\hbar^2$，再将 $\op{T}$ 的壳上矩阵元乘以 $-4\pi^2\mu/\hbar^2$ 换算为振幅，得到
    \begin{equation}
        \Im f^{(2)}(0)
        =\frac{4\pi^3\mu^2k}{\hbar^4}
            \int\abs{\mel{\bm k'}{\op{V}}{\bm k}}^2\d\Omega.
    \end{equation}
    另一方面，将一级振幅的模平方在全立体角上积分给出
    \[
        \sigma^{(2)}=\frac{16\pi^4\mu^2}{\hbar^4}
            \int\abs{\mel{\bm k'}{\op{V}}{\bm k}}^2\d\Omega .
    \]
    因而$\Im f^{(2)}(0)=k\sigma^{(2)}/(4\pi)$，
    正是在一致阶数上的光学定理。}

\reading{散射积分方程、Born近似与截面核对MIT 8.06讲义7\cite{mitscattering}。
    本章采用态和算符重写，并显式固定连续归一化。}
